\documentclass{extarticle}

% ========================================
% Page Layout and Geometry
% Font size and paper size are injected by the build wrapper via
% \PassOptionsToClass and \def\PAPERFORMAT before this file is \input.
% Defaults to a4 when PAPERFORMAT is undefined (e.g. direct compilation).
% ========================================
\ifdefined\PAPERFORMAT\else\def\PAPERFORMAT{a4}\fi
\def\fmtAfour{a4}%
\def\fmtAfive{a5}%
\def\fmtAsix{a6}%
\ifx\PAPERFORMAT\fmtAfour
\usepackage[a4paper, margin=0.75in, top=0.5in, bottom=0.75in]{geometry}
\else\ifx\PAPERFORMAT\fmtAfive
\usepackage[a5paper, margin=0.55in, top=0.4in, bottom=0.55in]{geometry}
\else
\usepackage[a6paper, margin=0.35in, top=0.3in, bottom=0.35in]{geometry}
\fi\fi

% ========================================
% Mathematical Packages
% ========================================
\usepackage{amsmath}
\usepackage{amssymb}
\usepackage{amsthm}
\usepackage{mathtools}

% ========================================
% Graphics and Color
% ========================================
\usepackage{graphicx}
\usepackage[dvipsnames]{xcolor}
\usepackage[stretch=10]{microtype}
\usepackage[htt]{hyphenat}

% ========================================
% Tables and Layout
% ========================================
\usepackage{multicol}
\usepackage{tabularx}
\usepackage{booktabs}
\usepackage{array}
\usepackage{adjustbox}
\usepackage{etoolbox}

% ========================================
% Lists and Enumeration
% ========================================
\usepackage{enumitem}
\setlist{itemsep=2pt, parsep=2pt}

% ========================================
% Code Listings
% ========================================
\usepackage{needspace}
\usepackage{listings}

% Rust language definition (not built into older listings versions)
\lstdefinelanguage{Rust}{
  keywords=[1]{as,async,await,break,const,continue,crate,dyn,else,enum,
    extern,false,fn,for,if,impl,in,let,loop,match,mod,move,mut,pub,ref,
    return,self,Self,static,struct,super,trait,true,type,union,unsafe,
  use,where,while},
  keywordstyle=[1]\color{blue!80!black}\bfseries,
  keywords=[2]{Option,Result,Some,None,Ok,Err,Box,Vec,String,
  HashMap,HashSet,Arc,Rc,Mutex,RwLock},
  keywordstyle=[2]\color{teal},
  keywords=[3]{i8,i16,i32,i64,i128,isize,u8,u16,u32,u64,u128,
  usize,f32,f64,bool,char,str},
  keywordstyle=[3]\color{violet},
  sensitive=true,
  comment=[l]{//},
  morecomment=[s]{/*}{*/},
  string=[b]",
  morestring=[b]',
}

\lstset{
  language=Rust,
  basicstyle=\small\ttfamily,
  breaklines=true,
  breakatwhitespace=false,
  keepspaces=true,
  showstringspaces=false,
  backgroundcolor=\color{gray!10},
  commentstyle=\color{ForestGreen}\itshape,
  keywordstyle=\color{blue!80!black}\bfseries,
  stringstyle=\color{red!70!black},
  numbers=none,
  % Left-bar frame: flows across page boundaries without open-ended boxes
  frame=leftline,
  framerule=3pt,
  rulecolor=\color{black!25},
  framexleftmargin=10pt,
  framexrightmargin=2pt,
  xleftmargin=0pt,
  xrightmargin=0pt,
  aboveskip=8pt,
  belowskip=8pt,
}

% ========================================
% Definition Boxes and Theorems
% ========================================
\usepackage[most,breakable]{tcolorbox}

% Definition box style
\newtcolorbox{definition}[1][]{
  breakable,
  colback=blue!5!white,
  colframe=blue!75!black,
  fonttitle=\bfseries,
  title=Definition,
  #1
}

% Concept box for important terms
\newtcolorbox{concept}[1]{
  breakable,
  colback=green!5!white,
  colframe=green!60!black,
  fonttitle=\bfseries,
  title=#1,
  before skip=10pt,
  after skip=10pt
}

% Important note box
\newtcolorbox{importantnote}{
  breakable,
  colback=orange!5!white,
  colframe=orange!75!black,
  fonttitle=\bfseries,
  title=Important Note,
  before skip=10pt,
  after skip=10pt
}

% Theorem environments
\newtheorem{theorem}{Theorem}[section]
\newtheorem{lemma}[theorem]{Lemma}
\newtheorem{proposition}[theorem]{Proposition}
\newtheorem{corollary}[theorem]{Corollary}

% ========================================
% Hyperlinks and References
% ========================================
\usepackage[
  colorlinks=true,
  linkcolor=blue!50!black,
  citecolor=red!50!black,
  urlcolor=blue!50!black,
  bookmarks=true,
  bookmarksopen=true,
  pdfauthor={Dr Jack Geraghty},
  pdftitle={Computational Audio and Image Analysis with the
  Spectrograms Library}
]{hyperref}

% Better equation referencing
\usepackage[capitalize,noabbrev]{cleveref}
\crefname{equation}{Equation}{Equations}
\crefname{section}{Section}{Sections}
\crefname{figure}{Figure}{Figures}
\crefname{table}{Table}{Tables}

% ========================================
% Bibliography
% ========================================
\usepackage[
  backend=biber,
  style=numeric,
  sorting=nyt,
  maxbibnames=99
]{biblatex}
\addbibresource{references.bib}

% ========================================
% Music Notation (if needed)
% ========================================
\usepackage{musicography}

% ========================================
% Custom Commands
% ========================================

% Math operators
\DeclareMathOperator{\DFT}{DFT}
\DeclareMathOperator{\IDFT}{IDFT}
\DeclareMathOperator{\FFT}{FFT}
\DeclareMathOperator{\IFFT}{IFFT}
\DeclareMathOperator{\STFT}{STFT}
\DeclareMathOperator{\CQT}{CQT}
\DeclareMathOperator{\MFCC}{MFCC}

% Complex conjugate
\newcommand{\conj}[1]{#1^*}

% Real and imaginary parts
\DeclareMathOperator{\Real}{Re}
\DeclareMathOperator{\Imag}{Im}

% Floor and ceiling
\DeclarePairedDelimiter{\floor}{\lfloor}{\rfloor}
\DeclarePairedDelimiter{\ceil}{\lceil}{\rceil}

% Absolute value and norm
\DeclarePairedDelimiter{\abs}{\lvert}{\rvert}
\DeclarePairedDelimiter{\norm}{\lVert}{\rVert}

% Code inline
\newcommand{\code}[1]{\texttt{#1}}

% Emphasis for technical terms on first use
\newcommand{\termdef}[1]{\textbf{\textit{#1}}}

% Rename the auto-generated listing list heading
\renewcommand{\lstlistlistingname}{List of Code Listings}
\renewcommand{\lstlistingname}{Listing}

% ========================================
% Typography and Spacing
% ========================================
\setlength{\parindent}{0pt}
\setlength{\parskip}{6pt}

% Better spacing in equations
\setlength{\jot}{8pt}

% Discourage orphan lines and widows
\widowpenalty=10000
\clubpenalty=10000
\displaywidowpenalty=10000

% Allow extra inter-word stretch as a last resort before declaring an
% overfull hbox — handles long \texttt and inline-math overflow.
\setlength{\emergencystretch}{2em}

% Keep section headings with their content: insert a page break if fewer
% than N lines remain on the page before the heading.
\preto
\section{\needspace{7\baselineskip}}
\preto
\subsection{\needspace{5\baselineskip}}
\preto
\subsubsection{\needspace{4\baselineskip}}

% On A5/A6, shrink table font so wide tables don't overflow the margins.
\ifx\PAPERFORMAT\fmtAfour\else
\AtBeginEnvironment{tabular}{\footnotesize}
\AtBeginEnvironment{tabularx}{\footnotesize}
\fi

% Kept for PDF metadata read by hyperref at \AtBeginDocument
\title{Computational Audio and Image Analysis with the Spectrograms Library}
\author{Dr Jack Geraghty}
\date{2026}

% ========================================
% Document
% ========================================
\begin{document}

% ----------------------------------------
% Title Page — nothing else lives here
% ----------------------------------------
\begin{titlepage}
  \centering
  \vspace*{5\baselineskip}
  \rule{\linewidth}{1pt}
  \vspace{0.6\baselineskip}

  {\Huge\bfseries Computational Audio and Image Analysis\\[0.4em]
  with the \texttt{Spectrograms} Library\par}

  \vspace{0.6\baselineskip}
  \rule{\linewidth}{1pt}

  \vspace{3\baselineskip}
  {\large Dr Jack Geraghty\par}

  \vfill
  {\normalsize Version 2.0 \textperiodcentered{} 2026\par}
\end{titlepage}

% ----------------------------------------
% Front Matter — roman numerals
% ----------------------------------------
\pagenumbering{roman}

\begin{abstract}
  This manual documents the mathematical foundations and computational
  implementation of the \texttt{spectrograms} library, a
  precision-generic Rust library for FFT-based audio and image
  analysis. It covers the Discrete Fourier Transform and fast FFT
  algorithms; windowed spectral analysis via the Short-Time Fourier
  Transform; perceptual frequency scales including the Mel scale,
  Equivalent Rectangular Bandwidth, and the Constant-Q Transform;
  advanced audio features such as Mel-frequency cepstral coefficients,
  chromagrams, and binaural spatial analysis via ITD, IPD, ILD, and
  ILR; two-dimensional Fourier methods for image filtering and edge
  detection; streaming convolution via the overlap-save method; and
  numerical precision considerations for single and double-precision
  computation. All algorithms are presented in closed-form mathematical
  notation alongside their Rust and Python implementations. The manual
  assumes familiarity with linear algebra, complex numbers, and basic
  discrete-time signal processing.
\end{abstract}
\clearpage

\tableofcontents
\clearpage

% These lists are only emitted if there are entries (flags written to
% the aux file at the end of each content pass and read on the next).
\makeatletter
\@ifundefined{@hasfigures}{}{\listoffigures\clearpage}
\@ifundefined{@hastables}{}{\listoftables\clearpage}
\@ifundefined{@haslistings}{}{\lstlistoflistings\clearpage}
\makeatother

% ----------------------------------------
% Main Matter — arabic numerals, page 1
% ----------------------------------------
\pagenumbering{arabic}

\input{manual.tex}

% ----------------------------------------
% Python API Reference Guide
% ----------------------------------------
\input{python_guide.tex}

% ----------------------------------------
% Write presence flags so the front matter knows what lists to show
% on the next compilation pass.
% ----------------------------------------
\makeatletter
\ifnum\value{figure}>0
\immediate\write\@mainaux{\string\gdef\string\@hasfigures{1}}
\fi
\ifnum\value{table}>0
\immediate\write\@mainaux{\string\gdef\string\@hastables{1}}
\fi
\ifnum\value{lstlisting}>0
\immediate\write\@mainaux{\string\gdef\string\@haslistings{1}}
\fi
\makeatother

% ----------------------------------------
% Bibliography
% ----------------------------------------
\clearpage
\printbibliography[title={References}]

\end{document}
